\section{Introduction}
%To detect piece movements, a camera is used, while a chess engine is responsible for generating the best possible moves the robot can take. The Gripper is controlled with a DC motor, an ADC and an H-bridge circuit. For moving and coordinating the robot motion, the operating system ROS2 framework is used in this project.

%The Git used for the project can be found in \cref{Git}, and the work distribution of the group can be seen in \cref{Work Distribution}.  

While online chess is easy and accessible, it lacks the physical presence of a real-world game. Therefore, a real-world chess game against a robot can be an attractive alternative.  

A robot will be set up to play a game of chess, autonomously against a human. The game should resemble a real-life chess game, so the player should have as little interaction with the robot as possible. The robot therefore needs to be able to play chess regardless of the position or rotation of the chessboard. The player should also be able to play as both white and black. 

For the robot to make its move, a gripper needs to be designed to pick up, move, and place chess pieces. A DC-motor is given as an actuator, therefore the gripper is to be designed to use a DC-motor. 

For version control, git and GitHub will be used and the repository for the project can be found in \cref{Git}, and the group's work distribution can be seen in \cref{Work Distribution}.